%% =============================================================================
%% Section 4 — Paradigm–SE Task Mapping
%% "How Do LLM-Based Agents Represent Code Repositories?"
%% ACM Transactions on Software Engineering and Methodology (TOSEM)
%% =============================================================================

\section{Paradigm--SE Task Mapping}
\label{sec:semapping}

This section answers RQ3 by mapping the seven repository representation
paradigms against seven canonical SE objectives: code completion (CC), code
generation from natural language (CG), automated program repair (APR), fault
and bug localization (FL), code search (CS), code translation (CT), and test
generation (TG).

\subsection{Mapping Methodology}

For each of the 159 primary studies, we recorded all SE objectives addressed.
A paradigm-task cell in Table~\ref{tab:paradigm-task} is marked \ding{51} if
at least three primary studies in that paradigm target that task, and
\textbf{(\ding{51})} if one or two studies do so.
An empty cell indicates no primary study in that paradigm addresses that task.

\begin{table*}[t]
  \caption{Paradigm--SE task coverage matrix.
           \ding{51}~= 3+ studies; \textbf{(\ding{51})}~= 1--2 studies;
           blank~= no studies.
           Numbers in parentheses indicate the count of primary studies.}
  \label{tab:paradigm-task}
  \small
  \begin{tabular}{lcccccccc}
    \toprule
    \textbf{Paradigm} & \textbf{Studies} &
    \textbf{CC} & \textbf{CG} & \textbf{APR} & \textbf{FL} &
    \textbf{CS} & \textbf{CT} & \textbf{TG} \\
    \midrule
    PAR-1 Dense Embedding     &  8 & \ding{51}    & \textbf{(\ding{51})} & & & \ding{51}    & & \\
    PAR-2 RAG                 & 35 & \ding{51}    & \ding{51} & \ding{51} & \textbf{(\ding{51})} & \ding{51} & & \ding{51} \\
    PAR-3 Graph               & 29 & \ding{51}    & \ding{51} & \textbf{(\ding{51})} & \ding{51} & \ding{51} & \ding{51} & \textbf{(\ding{51})} \\
    PAR-4 Execution           &  7 & & \textbf{(\ding{51})} & \ding{51} & \textbf{(\ding{51})} & & & \ding{51} \\
    PAR-5 Memory              &  4 & \textbf{(\ding{51})} & \ding{51} & \ding{51} & & & & \\
    PAR-6 Compression         & 10 & \ding{51}    & \ding{51} & & & & & \textbf{(\ding{51})} \\
    PAR-7 Agentic             & 41 & \ding{51}    & \ding{51} & \ding{51} & \ding{51} & \ding{51} & \ding{51} & \ding{51} \\
    \bottomrule
  \end{tabular}
\end{table*}

\subsection{Key Observations}

\paragraph{PAR-7 is the only universal paradigm.}
Agentic systems (PAR-7) appear in all seven SE task categories.
This reflects the generality of the agent-computer interface abstraction:
because agents can invoke arbitrary tools, they can be applied to any SE task
that can be reduced to a sequence of file operations.

\paragraph{PAR-1 is almost exclusively used for code search.}
Dense embeddings (PAR-1) dominate the code search task, where their ability to
map semantically similar code and natural language queries into a shared metric
space is directly applicable.
Their presence in code completion is more limited, typically as a retrieval
backbone inside a PAR-2 pipeline.

\paragraph{Automated program repair is the domain of PAR-4 and PAR-7.}
Execution-based representations (PAR-4) and agentic systems (PAR-7) together
account for the majority of APR studies.
This makes intuitive sense: repairing a bug requires understanding not just the
static code but its dynamic behaviour (PAR-4) or iteratively proposing and
testing patches (PAR-7).
PAR-2 appears in APR through retrieval of similar bugs and patches
\cite{wang2023rapgen}, but is not the dominant approach.

\paragraph{Graph-based representations (PAR-3) are the most task-diverse
non-agentic paradigm.}
PAR-3 appears in six of the seven task categories.
Its structural precision is especially valuable for fault localization
\cite{chen2025locagent} and code translation~\cite{ibrahimzada2025alphatrans},
tasks where explicit dependency information is essential.

\paragraph{Memory (PAR-5) is concentrated in generation and repair.}
The four PAR-5 studies target code generation and program repair, consistent
with the hypothesis that long-horizon tasks benefit most from experience-based
context~\cite{chen2025sweexp,pan2026prometheus}.

\paragraph{Compression (PAR-6) is absent from APR.}
No PAR-6 study addresses automated program repair.
This is consistent with the intuition that compression introduces information
loss that is particularly harmful for repair tasks, which require
\emph{precision} over large, complex code sections.
